\chapter{Relative quantum reduction across the discriminant}
\label{chap:vol2-relative-quantum-reduction}

Diagonalizing commuting matrices suggests that their invariant functions
should be functions of eigenvalue pairs. At repeated eigenvalues,
diagonalization no longer supplies a coordinate chart. The relation complex
also has odd functions that an ordinary quotient cannot see. A comparison
across this locus must therefore be a map of complexes, with its product
and coefficient ring specified.

In rank two the odd functions form a particularly small module. It has
rank one away from a quadratic cone's vertex, but needs two generators at
the vertex. We first compute this module and its restriction map.
A terminating differential operator then deforms the restriction over
$\mathbb C[\hbar]$. The result is a strict quasi-isomorphism, although it
has no strict algebra section. Its higher inverse makes this distinction
explicit. The final comparison with ordinary matrix reduction retains
the quantum parameter and trace rank at every matrix size.

\section{The relation module at a singular pair}
\label{sec:vol2-rank-two-koszul}

Work over $\mathbb C$ with cohomological differentials of degree one.
Identify the adjoint representation of $\mathfrak{sl}_2$ with the
three-dimensional representation of $G=SO_3(\mathbb C)$.
Normalize its invariant form and the equation generators so that the
commuting equation is the cross product. For even vectors $x,y$ and an
odd vector $\psi$, put
\[
 P=\mathbb C[x_1,x_2,x_3,y_1,y_2,y_3],
 \qquad
 K=P\otimes\Lambda(\psi_1,\psi_2,\psi_3),
 \qquad |\psi_i|=-1,
\]
\begin{equation}\label{eq:vol2-rank-two-koszul}
 dx=dy=0,\qquad d\psi=\mu,\qquad \mu=x\times y.
\end{equation}
The orientation is $\epsilon_{123}=1$. Exterior products are written
in the order $\psi_1,\psi_2,\psi_3$, and odd derivatives act on the left.
The relative complex is $K^G$, the invariants of the rational group
action. Taking these invariants imposes no positive gauge ghosts.

Define
\[
 a=x\cdot x,\qquad b=x\cdot y,\qquad c=y\cdot y,
 \qquad S=\mathbb C[a,b,c],\qquad \Delta=ac-b^2,
\]
and
\[
 \xi=x\cdot\psi,\qquad \eta=y\cdot\psi,\qquad z=\mu\cdot\psi.
\]
The elements $\xi,\eta$ are closed because $\mu$ is orthogonal to both
$x$ and $y$. To exhibit their relations, introduce
\[
 \begin{aligned}
 p&=x_1\psi_2\psi_3-x_2\psi_1\psi_3+x_3\psi_1\psi_2,\\
 q&=y_1\psi_2\psi_3-y_2\psi_1\psi_3+y_3\psi_1\psi_2,\\
 w&=\mu_1\psi_2\psi_3-\mu_2\psi_1\psi_3+\mu_3\psi_1\psi_2,
 \qquad r=\psi_1\psi_2\psi_3.
 \end{aligned}
\]
Here $p,q$ denote elements of exterior degree two, independently of the
plane coordinates used in Chapter~\ref{ch:vol2-derived-hamiltonian-actions}.
Their cohomological degrees are $-2$.

\begin{lemma}\label{lem:vol2-rank-two-primitives}
The differential satisfies
\begin{equation}\label{eq:vol2-rank-two-primitives}
 dz=\Delta,\qquad
 dp=b\xi-a\eta,\qquad dq=c\xi-b\eta,\qquad
 dr=w=\xi\eta,\qquad dw=0.
\end{equation}
\end{lemma}

\begin{proof}
Lagrange's identity gives $\mu\cdot\mu=ac-b^2$.
Identify a vector $v$ with the two-form
$v_1\psi_2\psi_3-v_2\psi_1\psi_3+v_3\psi_1\psi_2$.
Its differential is $(v\times\mu)\cdot\psi$.
The vector identities
\[
 x\times(x\times y)=bx-ay,\qquad
 y\times(x\times y)=cx-by
\]
give the middle two formulas. Expanding $d(\psi_1\psi_2\psi_3)$ gives
$w$, and expanding $(x\cdot\psi)(y\cdot\psi)$ gives the same two-form.
The last equation follows from $d^2r=0$.
\end{proof}

The invariant complex has eight coefficient generators. A polynomial
vector covariant is a polynomial map $v(x,y)$ satisfying
$v(gx,gy)=gv(x,y)$ for $g\in G$. The following calculation proves
the completeness of these generators, including on $\Delta=0$.

\begin{lemma}\label{lem:vol2-rank-two-covariants}
The scalar invariant ring is $S$. The polynomial vector covariants are
the free $S$-module generated by $x,y,\mu$. Consequently
\begin{equation}\label{eq:vol2-rank-two-components}
 \begin{aligned}
 (K^G)^0&=S,&
 (K^G)^{-1}&=S\xi\oplus S\eta\oplus Sz,\\
 (K^G)^{-2}&=Sp\oplus Sq\oplus Sw,&
 (K^G)^{-3}&=Sr.
 \end{aligned}
\end{equation}
\end{lemma}

\begin{proof}
Give a polynomial its separate degrees in $x$ and $y$.
The map $(x,y)\mapsto(a,b,c)$ is dominant: every nondegenerate
two-dimensional symmetric form embeds in the standard form on
$\mathbb C^3$. Thus $a,b,c$ are algebraically independent.
The determinant of the three columns $x,y,\mu$ is $\Delta$, a nonzero
polynomial. Hence these columns are linearly independent over the
fraction field and, in particular, over $S$.

It remains to prove that these explicit inclusions exhaust the
invariants. Let $V_\ell$ be the irreducible $SO_3$-representation of
spin $\ell$, of dimension $2\ell+1$. The elementary
$\mathfrak{sl}_2$ character formulas give
\[
 \operatorname{Sym}^m(V_1)
   =\bigoplus_{j=0}^{\lfloor m/2\rfloor}V_{m-2j},
 \qquad
 V_\ell\otimes V_k
   =\bigoplus_{n=|\ell-k|}^{\ell+k}V_n.
\]
For example, the first formula follows by comparing the coefficients
of $u^m$ in
\[
 \frac1{(1-uz)(1-u)(1-uz^{-1})}
 =\frac1{1-u^2}\sum_{\ell\geq0}u^\ell
     (z^\ell+z^{\ell-1}+\cdots+z^{-\ell}).
\]
The second follows by multiplying the two finite weight sums and
subtracting their successive highest weights.

The multiplicity of $V_0$ in $V_\ell\otimes V_k$ is one precisely
when $\ell=k$. The multiplicity of $V_1$ is one precisely when
$|\ell-k|\leq1\leq\ell+k$. Their bivariate generating series in the
polynomial ring are therefore
\[
 \frac1{(1-u^2)(1-uv)(1-v^2)},\qquad
 \frac{u+v+uv}{(1-u^2)(1-uv)(1-v^2)}.
\]
These are exactly the series of the already embedded ring $S$ and
module $Sx\oplus Sy\oplus S\mu$. Each bidegree is finite-dimensional,
so equality of dimensions makes those inclusions surjective.
The invariant form identifies vectors and covectors.
The invariant volume identifies their second exterior powers, and
the top exterior power is trivial. This gives
\eqref{eq:vol2-rank-two-components}.
\end{proof}

Set
\[
 B=S/(\Delta),\qquad
 M=(B\xi\oplus B\eta)/(b\xi-a\eta,\ c\xi-b\eta),
 \qquad |\xi|=|\eta|=-1.
\]
Let the involution on
$\mathbb C[s,t]\otimes\Lambda(\theta)$ change the signs of all three
generators, with $|\theta|=-1$. Restriction to a Cartan axis defines
\begin{equation}\label{eq:vol2-classical-rank-two-restriction}
 \rho:K^G\longrightarrow
     \bigl(\mathbb C[s,t]\otimes\Lambda(\theta)\bigr)^{\mathbb Z/2},
 \qquad
 x=(s,0,0),\quad y=(t,0,0),\quad\psi=(\theta,0,0).
\end{equation}
The target has zero differential.

\begin{theorem}\label{thm:vol2-classical-rank-two-restriction}
The map $\rho$ is a surjective quasi-isomorphism of commutative
differential graded algebras. It induces
\[
 H^0(K^G)=B,\qquad H^{-1}(K^G)=M,\qquad
 H^j(K^G)=0\quad(j\leq-2),
\]
with all products of negative-degree classes zero.
It sends $(a,b,c,\xi,\eta)$ to
$(s^2,st,t^2,s\theta,t\theta)$.
\end{theorem}

This is the traceless rank-two instance of the derived
Harish--Chandra comparison of Berest, Felder, Patotski, Ramadoss and
Willwacher (2015), Theorem~4.1 of~\cite{BFPRW2015}.
Here the coefficient presentation gives a direct proof.

\begin{proof}
The map on $Sz$ is multiplication by $\Delta$, hence injective.
On $Sp\oplus Sq$ the differential has matrix
\[
 \begin{pmatrix}b&c\\-a&-b\end{pmatrix}
\]
with determinant $\Delta$. It too is injective over the polynomial
domain $S$. Finally $Sr\to Sw$ is an isomorphism.
Thus the displayed four-term complex has no cohomology below
degree $-1$. Its degree-zero quotient is $B$, and its odd quotient is
$S^2/\langle(b,-a),(c,-b)\rangle$.
The two column relations imply $\Delta\xi=\Delta\eta=0$, so this
quotient is precisely $M$. The primitive $dr=\xi\eta$ proves
the product assertion.

The target's even part is
$\mathbb C[s^2,st,t^2]\cong B$. Its odd part consists of odd
polynomials in $s,t$, multiplied by $\theta$. It is generated by
$s\theta,t\theta$ over the even ring.
If even polynomials $F,G$ satisfy $Fs+Gt=0$, coprimality gives
$F=tH$ and $G=-sH$ with $H$ odd. Write $H=sP+tQ$ with $P,Q$
even. The relation is then
$P(st,-s^2)+Q(t^2,-st)$, exactly the relation module above.
Restriction is surjective on the displayed generators and induces
these cohomology isomorphisms.
\end{proof}

On $a\ne0$, the odd module is generated by $\xi$, since
$\eta=(b/a)\xi$. The same conclusion holds on $c\ne0$ with the
roles reversed. These two open sets cover the punctured cone.
At its vertex,
\begin{equation}\label{eq:vol2-rank-two-odd-vertex}
 M/(a,b,c)M=\mathbb C\xi\oplus\mathbb C\eta.
\end{equation}
This is the ordinary fibre of a cohomology module. It is not an
assertion that cohomology commutes with derived base change.
Replacing $M$ by a free line would already give the wrong ordinary
fibre. The restriction theorem uses no discriminant localization.

For $\mathfrak{gl}_2$, the scalar matrix coordinates supply a polynomial
factor $\mathbb C[u,v]$. The trace equation is identically zero,
so its relation generator $\zeta$, of degree $-1$, is closed.
Tensoring $\rho$ with
$\mathbb C[u,v]\otimes\Lambda(\zeta)$ gives the full rank-two
invariant-function comparison. Its degree $-2$ includes
$\mathbb C[u,v]\zeta M$. This extra odd direction cannot be discarded.
The statement concerns invariant functions, rather than an equivalence
of quotient stacks or of their module categories.

\section{The quantum differential and its coefficient operator}
\label{sec:vol2-relative-quantum-differential}

Use the matrix quantum parameter $\hbar$, independently of the
coefficient-plane parameter $\varepsilon$ in
Section~\ref{sec:vol2-coefficient-weyl-action}.
On polynomial symbols in the three canonical pairs $(x_i,y_i)$ set
\begin{equation}\label{eq:vol2-relative-moyal}
 F\star G
 =m\exp\left[
   \frac{\hbar}{2}\sum_{i=1}^3
   \left(\partial_{x_i}\otimes\partial_{y_i}
          -\partial_{y_i}\otimes\partial_{x_i}\right)
       \right](F\otimes G).
\end{equation}
Here $m$ is ordinary multiplication. Each product is polynomial in
$\hbar$. The $\psi_i$ still multiply
exteriorly. Define
\[
 \mathcal K_\hbar
   =\bigl(\mathbb C[x,y]\otimes\Lambda(\psi_1,\psi_2,\psi_3)
      \bigr)^G\otimes\mathbb C[\hbar]
\]
with this product and
\begin{equation}\label{eq:vol2-relative-quantum-differential}
 \delta_\hbar
 =\sum_i\mu_i\partial_{\psi_i}
  +\frac{\hbar^2}{4}
    \sum_{i,j,k}\epsilon_{ijk}
       \partial_{x_j}\partial_{y_k}\partial_{\psi_i}.
\end{equation}
The first summand uses ordinary multiplication of symbols.
The second summand is needed even though the moment functions
$\mu_i$ themselves are unchanged.

\begin{proposition}\label{prop:vol2-relative-quantum-dg}
The operator $\delta_\hbar$ is a square-zero derivation of
$\mathcal K_\hbar$ over $\mathbb C[\hbar]$.
\end{proposition}

\begin{proof}
Adjoin ghosts $\gamma^i$ of degree $1$ with
$\psi_i\gamma^j+\gamma^j\psi_i=\hbar\delta_i^j$.
The BRST element for the rotation moment is
\[
 \Theta=\sum_i\gamma^i\mu_i
  -\frac12\sum_{i,j,k}\epsilon_{ijk}\gamma^i\gamma^j\psi_k.
\]
The moment commutators are
$[\mu_i,\mu_j]_\star=\hbar\sum_k\epsilon_{ijk}\mu_k$.
Together with Jacobi they give $\Theta^2=0$, by the
Weyl--Clifford construction of
Section~\ref{sec:vol2-two-quantizations}.
Its polynomial inner derivation $[\Theta,-]/\hbar$ is therefore
square-zero.

Use symmetric symbols for the ghosts. On an expression independent
of $\gamma$, the terms that still contain $\gamma^i$ are its
infinitesimal simultaneous rotations, including the rotation of
$\psi$. They vanish on $G$-invariants. The remaining term is
\[
 \frac12\sum_i\left(
      \mu_i\star\partial_{\psi_i}F
        +\partial_{\psi_i}F\star\mu_i\right).
\]
Since $\mu_i$ is quadratic, the Moyal expansion stops at its second
derivative. The first derivatives cancel in the average, while
the second derivatives give exactly
\eqref{eq:vol2-relative-quantum-differential}.
The cubic ghost term has no higher contraction on a
$\gamma$-independent input. Thus this invariant subalgebra is
preserved by the inner derivation, proving both assertions.
\end{proof}

The presentation \eqref{eq:vol2-rank-two-components} remains a
presentation of polynomial symbols over $\mathbb C[\hbar]$.
It is a statement about coefficient modules, not about ordinary
multiplication being the quantum product. In particular
$\xi\star\eta=w$: the possible first contraction is proportional
to $\sum_i\psi_i^2$, hence vanishes.

Define operators on $S[\hbar]$ by
\begin{equation}\label{eq:vol2-rank-two-coefficient-operator}
 L=\partial_a\partial_c-\frac14\partial_b^2,\qquad
 T=1+\hbar^2L,\qquad
 T^{-1}F=\sum_{j\geq0}(-\hbar^2L)^jF.
\end{equation}
The inverse sum terminates because $L$ lowers polynomial degree
in $a,b,c$ by two. Neither $\hbar$ nor $\Delta$ is inverted.

\begin{lemma}\label{lem:vol2-rank-two-quantum-coefficients}
For every $F\in S[\hbar]$,
\begin{equation}\label{eq:vol2-rank-two-quantum-coefficients}
 \begin{aligned}
 \delta_\hbar(Fz)&=T(\Delta F),\\
 \delta_\hbar(Fp)&=T(bF)\xi-T(aF)\eta,\\
 \delta_\hbar(Fq)&=T(cF)\xi-T(bF)\eta,\\
 \delta_\hbar(Fr)&=T(F)w,\\
 \delta_\hbar(F\xi)&=\delta_\hbar(F\eta)
                  =\delta_\hbar(Fw)=0.
 \end{aligned}
\end{equation}
All products inside the arguments of $T$ are ordinary coefficient
products.
\end{lemma}

\begin{proof}
The chain rule gives
$\partial_{x_i}F=2x_iF_a+y_iF_b$ and
$\partial_{y_i}F=x_iF_b+2y_iF_c$. Therefore
\[
 \sum_{j,k}\epsilon_{ijk}
       \partial_{x_j}\partial_{y_k}F=4\mu_i LF.
\]
Applying the derivatives also to $\mu_i$ gives
\[
 \frac14\sum_{i,j,k}\epsilon_{ijk}
    \partial_{x_j}\partial_{y_k}(F\mu_i)
  =\Delta LF+(a\partial_a+b\partial_b+c\partial_c)F
       +\frac32F
  =L(\Delta F).
\]
This proves the first equation.
On $Fp$, the quantum coefficients of $\xi,\eta$ are
$bLF-F_b/2$ and $-aLF-F_c$.
On $Fq$, they are $cLF+F_a$ and $-bLF+F_b/2$.
These equal $L(bF),-L(aF),L(cF),-L(bF)$, respectively.
The top exterior differential has quantum correction $LF\,w$.
Antisymmetric contraction with $x$ or $y$ kills the corrections
on $F\xi,F\eta$.
Finally $F\mapsto T(F)$ is surjective, so every $Fw$ is the
differential of $T^{-1}(F)r$. Square-zero then gives its stated
differential.
\end{proof}

Let $U$ apply $T^{-1}$ to coefficients in
$S,S\xi,S\eta,Sw$ and act as the identity on
$Sz,Sp,Sq,Sr$. The coefficient formulas give
\begin{equation}\label{eq:vol2-relative-linear-comparison}
 U\delta_\hbar=\delta_0U.
\end{equation}
Thus $U$ is an integral chain isomorphism. Its inverse uses $T$
on the same four summands. This identity already computes the
underlying cohomology modules. To identify their product, the
restriction itself must be multiplicative.

\section{Strict restriction over the quantum ring}
\label{sec:vol2-strict-relative-restriction}

Let $W_\hbar(s,t)$ be the one-pair Weyl algebra in symmetric
symbols, so that $[s,t]_\star=\hbar$. Put
\[
 \mathcal H_\hbar
   =\bigl(W_\hbar(s,t)\otimes\Lambda(\theta)\bigr)^{\mathbb Z/2},
 \qquad |\theta|=-1,
\]
with zero differential and the simultaneous sign involution.
For a scalar symbol write $\rho_0F=F(s^2,st,t^2)$, where
$st$ in this substitution is the commuting symbol.
Define
\begin{equation}\label{eq:vol2-quantum-rank-two-restriction}
 \begin{aligned}
 \rho_\hbar(F)&=\rho_0(T^{-1}F),\\
 \rho_\hbar(F\xi+G\eta+Jz)
   &=\bigl(s\rho_0(T^{-1}F)+t\rho_0(T^{-1}G)\bigr)\theta,\\
 \rho_\hbar(\mathcal K_\hbar^{\leq-2})&=0.
 \end{aligned}
\end{equation}

\begin{theorem}\label{thm:vol2-quantum-rank-two-restriction}
The map $\rho_\hbar:\mathcal K_\hbar\to\mathcal H_\hbar$ is
a surjective strict quasi-isomorphism of differential graded
algebras over $\mathbb C[\hbar]$. Its kernel has a
$\mathbb C[\hbar]$-linear contraction. Consequently every base
change to a commutative $\mathbb C[\hbar]$-algebra preserves the quasi-isomorphism,
including every scalar specialization of $\hbar$.
\end{theorem}

\begin{proof}
As a chain map, $\rho_\hbar=\rho U$. The classical theorem
therefore proves surjectivity and the cohomology comparison.
We prove multiplicativity by the operators of multiplication
by the three invariant quadratics.

In $d$ canonical pairs, left multiplication on invariant scalar
symbols has the following operators:
\begin{align*}
 A_d={}&a+\hbar(a\partial_b+2b\partial_c)
   +\frac{\hbar^2}{4}
      (a\partial_b^2+4b\partial_b\partial_c
                    +4c\partial_c^2+2d\partial_c),\\
 B_d={}&b+\hbar(c\partial_c-a\partial_a)
   -\frac{\hbar^2}{4}
      (2a\partial_a\partial_b+b\partial_b^2
         +4b\partial_a\partial_c+2c\partial_b\partial_c
         +d\partial_b),\\
 C_d={}&c-\hbar(2b\partial_a+c\partial_b)
   +\frac{\hbar^2}{4}
      (4a\partial_a^2+4b\partial_a\partial_b
                       +c\partial_b^2+2d\partial_a).
\end{align*}
They follow by applying the two derivatives in the Moyal
product to $a,b,c$. For instance,
$a\star F=aF+\hbar x\cdot\partial_yF
 +\hbar^2\sum_i\partial_{y_i}^2F/4$,
and the chain rule gives $A_d$.

These differential operators satisfy
\begin{equation}\label{eq:vol2-rank-two-radial-intertwiners}
 A_3T=TA_1,\qquad B_3T=TB_1,\qquad C_3T=TC_1.
\end{equation}
Here is the coefficient verification. The commutators of $L$
with multiplication by $a,b,c$ are
$\partial_c,-\partial_b/2,\partial_a$.
It commutes with the three first-order coefficients of $\hbar$.
Its commutators with the coefficients of $\hbar^2$ in
$A_1,B_1,C_1$ are, respectively,
$L\partial_c,-L\partial_b/2,L\partial_a$.
Thus $[T,A_1]=\hbar^2\partial_cT=(A_3-A_1)T$,
and the other two equations follow with the displayed constants.

Evaluation by $\rho_0$ changes $A_1,B_1,C_1$ into actual
one-pair left multiplication by $s^2,st,t^2$.
Equation~\eqref{eq:vol2-rank-two-radial-intertwiners}
therefore proves that $\rho_\hbar$ intertwines multiplication
by $a,b,c$ on the scalar sector.
These three quadratics generate its invariant Weyl algebra:
a star monomial has its ordinary monomial as leading symbol,
and subtraction lowers polynomial degree. Induction proves
generation over $\mathbb C[\hbar]$.

On the coefficient column $(F,G)$ of $F\xi+G\eta$,
the corresponding left multiplication operators are
\[
 \begin{pmatrix}
 A_d&\hbar+\frac{\hbar^2}{2}\partial_b\\
 0&A_d+\hbar^2\partial_c
 \end{pmatrix},
 \qquad
 \begin{pmatrix}
 B_d-\frac\hbar2-\frac{\hbar^2}{4}\partial_b
    &-\frac{\hbar^2}{2}\partial_a\\
 -\frac{\hbar^2}{2}\partial_c
    &B_d+\frac\hbar2-\frac{\hbar^2}{4}\partial_b
 \end{pmatrix},
\]
\[
 \begin{pmatrix}
 C_d+\hbar^2\partial_a&0\\
 -\hbar+\frac{\hbar^2}{2}\partial_b&C_d
 \end{pmatrix}.
\]
For example, multiplying the coefficient $G$ and the single
$y$ in $\eta$ accounts for the two off-diagonal terms of the
first matrix. The other entries follow from the same product
rule. Their extra entries have constant coefficients, so they
commute with $T$. The scalar intertwiners prove the column
intertwiners. Replacing $\hbar$ by $-\hbar$ gives right
multiplication as well.

Give $x,y$ weight one and $\psi,\hbar$ weight two.
The scalar and $Sz$ sectors have even weight, while
$S\xi,S\eta$ have odd weight. Scalar star multiplication
therefore preserves the killed $Sz$ sector.
Products of total cohomological degree at most $-2$ have
zero image, and two odd target factors have product zero
because $\theta^2=0$. This exhausts the products and proves
strict multiplicativity.

Choose a contraction of the classical kernel in each finite
weight. Such a contraction exists by the classical
quasi-isomorphism: split each vector space into boundaries
and chosen preimages, and send boundaries back to their
preimages. Extend it $\mathbb C[\hbar]$-linearly and conjugate
by $U$. This contracts the quantum kernel.
Together with a chain section of the surjection, it gives a
contraction of the whole comparison. Tensoring these
identities with any commutative coefficient algebra proves the base-change
assertion without a flatness assumption.
\end{proof}

The correction is already visible in three small values:
\begin{equation}\label{eq:vol2-relative-first-quantum-corrections}
 \rho_\hbar(ac)=s^2t^2-\hbar^2,\qquad
 \rho_\hbar(b^2)=s^2t^2+\frac12\hbar^2,\qquad
 \delta_\hbar z=\Delta+\frac32\hbar^2.
\end{equation}
The naive substitution $\rho_0$ would send the last boundary
to $3\hbar^2/2$, so it is not the quantum chain map.
The strict comparison instead identifies the even cohomology
relation
\[
 \frac12(a\star c+c\star a)-b\star b=\frac34\hbar^2.
\]
Its odd relations retain their primitives:
\[
 (b+\tfrac32\hbar)\star\xi-a\star\eta=\delta_\hbar p,
 \qquad
 c\star\xi-(b-\tfrac32\hbar)\star\eta=\delta_\hbar q.
\]
Tensoring with the independent scalar Weyl pair and its closed
central antighost gives the relative $\mathfrak{gl}_2$ comparison.
The absolute complex still has its positive ghosts.

\section{An inverse with higher components}
\label{sec:vol2-relative-higher-inverse}

A strict projection need not have a strict section. To construct
an inverse of $\rho_\hbar$ with its homotopies, choose the classical
section on the Cartan monomial basis by
\[
 i_0(s^mt^n)=
 \begin{cases}
  a^{m/2}c^{n/2},&m,n\text{ even},\\
  a^{(m-1)/2}bc^{(n-1)/2},&m,n\text{ odd},
 \end{cases}
\]
\[
 i_0(s^mt^n\theta)=
 \begin{cases}
  a^{(m-1)/2}c^{n/2}\xi,&m\text{ odd},\ n\text{ even},\\
  a^{m/2}c^{(n-1)/2}\eta,&m\text{ even},\ n\text{ odd}.
 \end{cases}
\]
In each finite weight split the classical complex into these
cohomology representatives, boundaries and chosen preimages of
boundaries. Let $H_0$ send a boundary to its preimage and vanish on
the other two summands. Choose $H_0\Delta=z$ and $H_0w=r$.
Then
\[
 \delta_0H_0+H_0\delta_0=1-i_0\rho,\qquad
 H_0^2=H_0i_0=\rho H_0=0.
\]
Extend these maps $\mathbb C[\hbar]$-linearly and put
\[
 i_\hbar=U^{-1}i_0,\qquad
 H_\hbar=U^{-1}H_0U.
\]
They give the same normalized contraction identities for
$(\rho_\hbar,i_\hbar,H_\hbar)$.

Let $s$ lower cohomological degree by one.
On the tensor coalgebras of the suspended underlying nonunital
algebras use the codifferentials determined by
\[
 b_1(sa)=-s\delta_\hbar a,\qquad
 b_2(sa,sb)=(-1)^{|a|}s(a\star b).
\]
The target algebra has only $b_2$.
These coalgebras have the deconcatenation coproduct; their maps
are determined by their components into tensors of length one.
No scalar augmentation of a Weyl algebra is required.
Put $H_s=-sH_\hbar s^{-1}$ and $I_1=si_\hbar s^{-1}$.
Given the components below arity $n$, set
\begin{equation}\label{eq:vol2-relative-inverse-recursion}
 O_n=(b_{\mathcal K}I_{<n}-I_{<n}b_{\mathcal H})_n,
 \qquad I_n=-H_sO_n.
\end{equation}
Here $I_{<n}$ denotes the coalgebra map with the prescribed
lower components and zero arity-$n$ component.

\begin{theorem}\label{thm:vol2-relative-higher-inverse}
This recursion defines a strictly unital $A_\infty$ morphism
$I:\mathcal H_\hbar\to\mathcal K_\hbar$ whose first component
is a quasi-isomorphism and for which $\rho_\hbar I$ is the
strict identity. Its only possible components are $I_1,I_2,I_4$.
For the commuting target symbol $B=st$,
\begin{equation}\label{eq:vol2-relative-fourth-inverse}
 I_2(sB,sB)=-sz,\qquad
 I_4(sB,sB,sB,sB)=-\hbar sr.
\end{equation}
\end{theorem}

\begin{proof}
The equation $b^2=0$, together with the lower morphism equations,
makes $O_n$ closed for the unary differential.
Strictness of $\rho_\hbar$ and
$\rho_\hbar i_\hbar=1$, $\rho_\hbar H_\hbar=0$
put $O_n$ in the kernel of the suspended projection.
The contraction identity makes $b_1(-H_sO_n)=-O_n$.
Thus \eqref{eq:vol2-relative-inverse-recursion} solves the next
morphism equation. Every map is polynomial on its input.

The special vanishing follows from the weights and the four
cohomological degrees. On inputs with $k$ odd factors, the
unsuspended third residual has degree $-1-k$ and weight parity
$k$. For $k=0,1,2$ it lies in $S[\hbar]z$,
$S[\hbar]p\oplus S[\hbar]q$, and $S[\hbar]r$,
respectively. The differential is injective on each of these
sectors, by the coefficient formulas and the determinant
$\Delta$ of the middle matrix. Closedness forces the residual
to vanish. For larger $k$ its degree is absent. Hence $I_3=0$.

The fourth component has unsuspended degree $-3-k$, so it
can occur only on four even inputs and takes values in
$S[\hbar]r$. All components of arity at least five would
have degree below $-3$. The fifth residual on even inputs
lies in $S[\hbar]r$ and is closed, hence zero by injectivity.
Every later residual has absent degree. This proves all
morphism equations, including those beyond the last nonzero
component. Since the section preserves the unit and the
homotopy annihilates its image, the normalized recursion
vanishes on a unit input in every higher arity.

For the stated example,
\[
 b\star b=b^2-\frac34\hbar^2,\qquad
 i_\hbar(B\star B)=ac+\frac34\hbar^2.
\]
Their difference is $-\delta_\hbar z$, giving $I_2=-sz$.
The moment commutator and exterior signs give
\[
 z\star z=\hbar w,\qquad \delta_\hbar r=w.
\]
Consequently the fourth residual is
$b_2(-sz,-sz)=-\hbar sw$, and the recursion gives
$I_4=-\hbar sr$.
\end{proof}

For even inputs $F,G$, the binary component is particularly
explicit. Set
\[
 D=i_\hbar(F)\star i_\hbar(G)-i_\hbar(F\star G).
\]
Then $T^{-1}D$ is divisible by $\Delta$ in the polynomial
ring. The unsuspended binary value is
$(T^{-1}D/\Delta)z$. This is exact polynomial division,
not localization at the discriminant.
The fourth value is obtained by multiplying two binary
values and applying $-T^{-1}$ to the resulting $w$ coefficient.
The target still has its ordinary binary product and no
higher multiplications.

\section{Changing a section}
\label{sec:vol2-relative-sections}

The nonzero fourth value belongs to a chosen inverse.
There are two separate questions: whether the projection has
a strict section on the whole algebra, and whether a
particular higher value survives a change of section.

\begin{proposition}\label{prop:vol2-relative-no-strict-section}
The map $\rho_\hbar$ has no $\mathbb C[\hbar]$-linear
dg-algebra section.
\end{proposition}

\begin{proof}
A section would specialize at $\hbar=0$ to an algebra
section of $S\to S/(\Delta)$ in degree zero.
Its lifts of $a,b,c$ would have the form
$A=a+\Delta\alpha$, $B=b+\Delta\beta$,
$C=c+\Delta\gamma$. The relation $AC=B^2$ would imply
\[
 0=\Delta\left(
  1+a\gamma+c\alpha-2b\beta
       +\Delta(\alpha\gamma-\beta^2)\right).
\]
Since $S$ is a domain, the parenthesis would be zero.
Its constant term is one, a contradiction.
\end{proof}

\begin{proposition}\label{prop:vol2-relative-polynomial-section}
There is a unital weight-preserving chain section $i'$ of
$\rho_\hbar$ whose restriction to the polynomial star-subalgebra
$\mathbb C[\hbar][B]$, $B=st$, is the algebra map
\[
 i'(B^{\star n})=b^{\star n}.
\]
A normalized contraction with this section has a recursive
$A_\infty$ inverse whose higher components all vanish when
every input belongs to this subalgebra.
\end{proposition}

\begin{proof}
The leading symbols $(st)^n$ show that the star powers
$B^{\star n}$ are linearly independent over $\mathbb C[\hbar]$.
Multiplicativity of $\rho_\hbar$ makes the proposed lift a
section on their polynomial algebra.
An even basis of the entire target is given by
\[
 A^{\star i}\star B^{\star j}\ (i\geq1),\qquad
 B^{\star j},\qquad
 B^{\star j}\star C^{\star k}\ (k\geq1),
 \qquad A=s^2,\quad C=t^2.
\]
Their leading symbols partition all even monomials in $s,t$.
The lower terms are triangular by polynomial degree, so
these star words are a basis over the parameter ring.
Lift them by the same star words in $a,b,c$.
Choose closed lifts of the odd monomial basis using the
surjective restriction on $S\xi\oplus S\eta$.
This gives a unital weight-preserving chain section $i'$.

From the previous normalized contraction put
$H'=H_\hbar(1-i'\rho_\hbar)$.
The chain-section identity gives
\[
 \delta_\hbar H'+H'\delta_\hbar=1-i'\rho_\hbar.
\]
Moreover $H'i'=\rho_\hbar H'=0$.
The identities $\rho_\hbar H_\hbar=H_\hbar^2=0$
give $(H')^2=0$.
The binary defect on the polynomial subalgebra is zero
because $i'$ is multiplicative there. Inductively every
higher residual consists of lower higher components on
that same subalgebra, so the recursion chooses zero.
\end{proof}

The change already appears on the square:
\[
 i_\hbar(B\star B)=ac+\frac34\hbar^2,\qquad
 i'(B\star B)=b^2-\frac34\hbar^2.
\]
Their difference is $-\delta_\hbar z$.
The target product has not changed. In particular the
repeated-$B$ fourth value in
\eqref{eq:vol2-relative-fourth-inverse} can be removed.

Higher data must nevertheless occur somewhere for the
weight-preserving recursive inverses. A classical section
cannot be multiplicative by
Proposition~\ref{prop:vol2-relative-no-strict-section}, so some
even pair $F,G$ has binary defect $\Delta\phi_0$ with
$\phi_0\ne0$. Its quantum binary primitive has the form
$\phi_\hbar z$, where $\phi_\hbar\bmod\hbar=\phi_0$.
Invariant coefficients Poisson-commute with each $\mu_i$.
It follows that
\[
 (\phi_\hbar z)\star(\phi_\hbar z)
   =\bigl(\hbar\phi_0^2+O(\hbar^2)\bigr)w\ne0.
\]
Here $O(\hbar^2)$ denotes a polynomial divisible by $\hbar^2$.
The degree and parity proof above still makes $I_3=0$.
On $(F,G,F,G)$ the fourth residual is this nonzero product,
so injectivity of $\delta_\hbar:S[\hbar]r\to S[\hbar]w$
forces a nonzero fourth value.

These facts are compatible with formality.
A dg algebra is \emph{formal} if a sequence of dg-algebra
quasi-isomorphisms identifies it with its cohomology algebra
with zero differential.
The strict map $\rho_\hbar$ already proves that
$\mathcal K_\hbar$ is formal. Failure of a section in its
reverse direction does not contradict that conclusion.
Neither a chosen fourth value nor its physical interpretation
is an invariant obstruction merely by being nonzero.

\section{The odd coefficient bimodule}
\label{sec:vol2-relative-odd-bimodule}

The factor $\theta$ makes the product of two odd cohomology
classes zero. Removing that factor gives a different
operation worth retaining. Write
\[
 A_\hbar=W_\hbar(s,t)^{\mathrm{even}},\qquad
 M_\hbar=W_\hbar(s,t)^{\mathrm{odd}}.
\]
Parity here refers to the sign involution on the even
coordinates, not to cohomological degree. Multiplication
gives a bimodule map
$M_\hbar\otimes_{A_\hbar}M_\hbar\to A_\hbar$.

\begin{proposition}\label{prop:vol2-relative-odd-bimodule}
Over $\mathbb C[\hbar]$, the image of this multiplication
is the ideal $(s^2,st,t^2,\hbar)$ in $A_\hbar$, with
cokernel $\mathbb C$ at $s=t=\hbar=0$.
After inverting $\hbar$, multiplication is an isomorphism,
with inverse
\[
 a\longmapsto
   \hbar^{-1}(s\otimes t-t\otimes s)a.
\]
\end{proposition}

\begin{proof}
The four basic products $s\star s,s\star t,t\star s,t\star t$
generate $s^2,st,t^2,\hbar$. Every even ordered monomial of
positive degree splits into two odd monomials.
Their products span all the other even monomials modulo
lower terms, by induction on degree. Thus their image is
the stated ideal. Its quotient leaves only the scalar origin.

In the localized ring, multiplication of the displayed
inverse gives $a$ because $s\star t-t\star s=\hbar$.
For the reverse composite on $m\otimes n$, use that
$t\star m$ and $s\star m$ are even. Balancing them across
the tensor product gives numerator
\[
 s\otimes(t\star m)\star n
   -t\otimes(s\star m)\star n
  =(s\star t-t\star s)\star m\otimes n
  =\hbar m\otimes n.
\]
This proves the second identity.
\end{proof}

Only this last invertibility assertion inverts the parameter.
The strict relative comparison was integral.
An interpretation of the coefficient bimodule product as
composition of physical lines requires a realization map;
the algebraic calculation alone does not construct one.

\section{The ordinary reduction at every matrix size}
\label{sec:vol2-integral-matrix-reduction}

The rank-two complex has retained the relation classes.
There is also an ordinary quantum reduction at every matrix
size, defined across the discriminant. Its global trace
action is an input to a derived comparison, rather than
a replacement for that comparison.

Fix $N\geq1$ and set $R=\mathbb C[\hbar]$.
Let $W_{N,\hbar}$ be generated by two $N$ by $N$ matrices
$X,Y$ with
\[
 [X_{ij},X_{kl}]=[Y_{ij},Y_{kl}]=0,\qquad
 [X_{ij},Y_{kl}]=\hbar\delta_{il}\delta_{jk}.
\]
The group $GL_N(\mathbb C)$ acts by simultaneous conjugation.
Use the ordered moment matrix from
Section~\ref{sec:vol2-two-quantizations},
\[
 \widehat\mu=XY-YX-N\hbar I_N,
 \qquad
 I_{N,\hbar}=\sum_{i,j}W_{N,\hbar}\widehat\mu_{ij}.
\]
This is a left ideal. Its invariant intersection is
two-sided in the invariant algebra, since invariant
operators commute with the infinitesimal conjugation
operators. Define the ordinary reduction
\begin{equation}\label{eq:vol2-ordinary-matrix-reduction}
 \mathcal A_{N,\hbar}
 =W_{N,\hbar}^{GL_N}/
       (I_{N,\hbar}\cap W_{N,\hbar}^{GL_N}).
\end{equation}
Replacing the full left ideal by the two-sided ideal it
generates would define a different problem.

Let $\mathfrak g=\mathfrak{gl}_N$ and let $D(\mathfrak g)$
be the polynomial differential operators.
Its order filtration $F_m$ allows arbitrary polynomial
coefficients and at most $m$ derivatives, with $F_{-1}=0$.
The Rees algebra is
\[
 \operatorname{Rees}_F D(\mathfrak g)
  =\bigoplus_{m\geq0}\hbar^mF_mD(\mathfrak g)
       \subset D(\mathfrak g)[\hbar].
\]
The assignments
$X_{ij}\mapsto x_{ij}$ and
$Y_{ij}\mapsto-\hbar\partial_{x_{ji}}$ identify it with
$W_{N,\hbar}$. The Weyl relations span the monomials with
all $X$ before all $Y$. Their differential-operator images
are independent after inverting $\hbar$, so they have no
nonzero $R$-linear relation before inversion either.
They span the Rees algebra together with powers of $\hbar$.
Thus $X$ has Rees order zero, while $Y$ and $\hbar$ have
order one. These are not cohomological degrees.

Let $\tau(a)$ be the differentiated pullback conjugation
action, normalized by $\tau(a)X=[X,a]$ on coordinate
functions. The Weyl relations give
\[
 \operatorname{Tr}(a\widehat\mu)=\hbar\tau(a).
\]
In particular $\operatorname{Tr}\widehat\mu=0$.
Put $J_D=(D(\mathfrak g)\tau(\mathfrak g))^{GL_N}$.
Let $J_{\mathrm{cl}}$ be the ideal generated by the entries
of the commutator of a pair in $\mathfrak g\oplus\mathfrak g$.
The trace pairing identifies the second factor with
cotangent momentum.

Write $\mathfrak t$ for diagonal matrices and $S_N$ for
its Weyl group. On the distinct-eigenvalue locus, the raw
radial operator $\operatorname{rad}(D)$ acts on restrictions
of invariant functions. Its conjugated form is
\[
 \Phi(D)=\delta\,\operatorname{rad}(D)\,\delta^{-1},
 \qquad
 \delta=\prod_{i<j}(\lambda_i-\lambda_j).
\]
The radial definition begins on this open set.
The following theorem is the extension statement that
allows it to be used across the discriminant.

\begin{theorem}[Filtered Harish--Chandra comparison]
\label{thm:vol2-filtered-harish-chandra}
For $\mathfrak g=\mathfrak{gl}_N$, the conjugated radial map
induces an algebra isomorphism
\[
 D(\mathfrak g)^{GL_N}/J_D
   \xrightarrow{\ \sim\ }D(\mathfrak t)^{S_N}.
\]
It is strict for the order filtrations. Its associated
graded map is restriction of commuting pairs to diagonal
pairs, and
\[
 \operatorname{gr}_FJ_D
       =(\sqrt{J_{\mathrm{cl}}})^{GL_N}
       =J_{\mathrm{cl}}^{GL_N}.
\]
\end{theorem}

The algebra isomorphism and associated graded assertions are
the Harish--Chandra results recorded by Etingof and Ginzburg
(2001), Theorems~6.8--6.10 of~\cite{EG2001HC}.
The equality of the invariant ideal with its radical is
Gan--Ginzburg (2005), Theorem~1.2.1 and
Section~2.8 of~\cite{GG2005Almost}.
Strictness follows by lifting a leading symbol and subtracting:
the differential order decreases at each step.
The equality of invariant ideals does not assume that the
whole commuting scheme is reduced.

\begin{lemma}\label{lem:vol2-strict-moment-orders}
Every $D\in J_D\cap F_mD(\mathfrak g)$ is a finite invariant
sum of terms $P_a\tau(a)$ with
$P_a\in F_{m-1}D(\mathfrak g)$.
For $m=0$ the only such $D$ is zero.
\end{lemma}

\begin{proof}
The leading symbol $\sigma_m(D)$ belongs to
$J_{\mathrm{cl}}^{GL_N}$.
Each moment generator has momentum degree one.
Its homogeneous ideal expression can therefore be written
\[
 \sigma_m(D)=\sum_a p_a\,\sigma_1(\tau(a)),
 \qquad \deg_{\mathrm{momentum}}p_a=m-1.
\]
Only finitely many coefficients occur.
The map from coefficient vectors to the ideal is equivariant.
Rational $GL_N$-invariants are exact in characteristic zero:
on a finite-dimensional rational module, average over the
unitary compact form. Applying this to a degree-bounded
module containing the coefficients gives an invariant
coefficient vector with the same image.

Normal quantization, placing coordinates before derivatives,
is equivariant under the linear action. It lifts the $p_a$
to $P_a$ of order at most $m-1$.
The invariant operator $\sum_aP_a\tau(a)$ has the same
leading symbol as $D$. Subtract it and repeat at smaller
order. This process terminates.
The degree-zero momentum part of $J_{\mathrm{cl}}$ is zero,
so the last order-zero remainder vanishes.
\end{proof}

An ideal $J$ is $\hbar$-saturated if $\hbar u\in J$
implies $u\in J$. Let $W^{\mathrm{Cartan}}_{N,\hbar}$
be the Weyl algebra on $N$ pairs $\lambda_i,\pi_i$ with
$[\lambda_i,\pi_j]=\hbar\delta_{ij}$ and all other
generator commutators zero.

\begin{theorem}\label{thm:vol2-integral-harish-chandra}
The conjugated radial map has an $R$-algebra realization
\[
 \Phi_{N,\hbar}:\mathcal A_{N,\hbar}
   \xrightarrow{\ \sim\ }
   (W^{\mathrm{Cartan}}_{N,\hbar})^{S_N}.
\]
Moreover
\[
 I_{N,\hbar}\cap W_{N,\hbar}^{GL_N}
     =\operatorname{Rees}_FJ_D.
\]
This invariant ideal is $\hbar$-saturated, the quotient
is $R$-free, and its special fibre is
\[
 \mathcal A_{N,\hbar}/\hbar\mathcal A_{N,\hbar}
   =\mathbb C[\lambda_1,\pi_1,\ldots,\lambda_N,\pi_N]^{S_N}.
\]
\end{theorem}

\begin{proof}
The conjugation action preserves every filtration step, so
invariants give
$W_{N,\hbar}^{GL_N}
 =\operatorname{Rees}_F(D(\mathfrak g)^{GL_N})$.
In each Rees degree an element of the left moment ideal
has coefficient in $D(\mathfrak g)\tau(\mathfrak g)$.
Its invariant part therefore belongs to
$\operatorname{Rees}_FJ_D$.

Conversely, let $\hbar^mD$ have $D\in F_mJ_D$.
The preceding lemma expresses
$D=\sum_aP_a\tau(a)$ with $P_a\in F_{m-1}D(\mathfrak g)$.
Hence
\[
 \hbar^mD=\sum_a(\hbar^{m-1}P_a)(\hbar\tau(a))
       \in I_{N,\hbar}.
\]
This element is invariant, proving equality of the ideals.
The strict filtered Harish--Chandra isomorphism now
identifies the quotient with the Rees algebra of
$D(\mathfrak t)^{S_N}$.
Its Weyl generators are the displayed $\lambda_i,\pi_i$.

Weyl symmetrization identifies its underlying $R$-module
equivariantly with the ordinary polynomial module.
A complex basis of the $S_N$-invariant polynomial space
therefore gives an $R$-basis of the quotient.
It has no $\hbar$-torsion. Thus if $\hbar u$ maps to zero,
$u$ maps to zero too, which is saturation.
The Rees special fibre is its associated graded algebra,
giving the final formula.
\end{proof}

Saturation here concerns the invariant intersection.
No assertion about saturation of the entire left ideal
was required. At $\hbar=0$, $\pi_i$ remains a momentum
symbol in the intrinsic special fibre. Substituting
$\hbar=0$ into the external representation
$\pi_i=-\hbar\partial_{\lambda_i}$ would instead annihilate
it. Specializing that representation does not compute
the Rees special fibre.

\section{The global ordinary trace action}
\label{sec:vol2-global-ordinary-moyal-action}

For $f,g\in R[q,p]$, let
$B_\hbar(f,g)=(f\star_\hbar g-g\star_\hbar f)/\hbar$
be the scalar Moyal bracket. Its numerator is polynomial
and divisible by $\hbar$.
For a monomial define the matrix-word symmetrized trace
\[
 J_N(q^ap^b)
  =\binom{a+b}{a}^{-1}\sum_{v}\operatorname{Tr}v(X,Y),
\]
where $v$ runs over distinct ordered products with $a$
occurrences of $X$ and $b$ of $Y$.
Extend $J_N$ $R$-linearly. In particular $J_N(1)=N$.
Let $\operatorname{Op}_W(f)(\lambda,\pi)$ denote
one-pair Weyl symmetrization.

\begin{theorem}\label{thm:vol2-global-moyal-trace}
For every matrix size $N$ and every polynomial $f$,
\begin{equation}\label{eq:vol2-global-moyal-trace-image}
 \Phi_{N,\hbar}([J_N(f)])
     =\sum_{i=1}^N\operatorname{Op}_W(f)(\lambda_i,\pi_i).
\end{equation}
Consequently
\[
 \frac{[J_N(f),J_N(g)]}{\hbar}-J_N(B_\hbar(f,g))
       \in I_{N,\hbar}\cap W_{N,\hbar}^{GL_N}.
\]
The induced map
$R[q,p]/R1\to\mathcal A_{N,\hbar}/R1$ is injective
and preserves the normalized Lie brackets.
\end{theorem}

\begin{proof}
Multiplication by $\operatorname{Tr}X^d$ restricts to
$\sum_i\lambda_i^d$.
The radial image of the quadratic momentum trace is
\[
 \operatorname{rad}(\operatorname{Tr}Y^2)
 =\hbar^2\left(
   \sum_i\partial_{\lambda_i}^2
    +2\sum_{i<j}
       \frac{\partial_{\lambda_i}-\partial_{\lambda_j}}
            {\lambda_i-\lambda_j}\right).
\]
To see the mixed derivative coefficient, perturb entries
$(i,j)$ and $(j,i)$ by $u,v$. Their two eigenvalues satisfy
$(z-\lambda_i)(z-\lambda_j)-uv=0$.
Their mixed second derivatives are
$\pm(\lambda_i-\lambda_j)^{-1}$, yielding the displayed
first-order terms.
The Vandermonde polynomial $\delta$ is harmonic:
its Laplacian is alternating of lower degree, hence
divisible by $\delta$ and therefore zero.
Conjugation by $\delta$ removes the first-order terms.
Thus $\Phi_{N,\hbar}([\operatorname{Tr}Y^2])
 =\sum_i\pi_i^2$.

Put $H=-\operatorname{Tr}Y^2/2$.
The polynomial derivation $D=[H,-]/\hbar$ satisfies
$DX=Y$ and $DY=0$.
Its exponential terminates on any polynomial in the
matrix entries and gives
\[
 e^{tD}\operatorname{Tr}X^d
      =\operatorname{Tr}(X+tY)^d.
\]
Multiplicativity of $\Phi_{N,\hbar}$ and torsion-freeness
of its target show that it intertwines $D$ with
$\lambda_i\mapsto\pi_i$, $\pi_i\mapsto0$.
The coefficient of $t^b$, with $a+b=d$, is the sum of
the matrix products defining $J_N(q^ap^b)$.
The same coefficient on the Cartan side is the sum of
the corresponding one-pair Weyl products.
Dividing by the same nonzero binomial coefficient proves
\eqref{eq:vol2-global-moyal-trace-image}.

Different Cartan pairs commute, and one-pair Weyl
quantization intertwines Moyal products. The stated
bracket defect has zero radial image.
The exact kernel in
Theorem~\ref{thm:vol2-integral-harish-chandra} puts it
in the invariant moment ideal.
The division by $\hbar$ is the polynomial quotient of
a divisible numerator throughout.

Finally suppose $[J_N(f)]$ is scalar.
Taking Cartan Weyl symbols gives
$\sum_i f(\lambda_i,\pi_i)=\alpha$ for some $\alpha\in R$.
Set all pairs except the first equal to zero. Then
$f(q,p)+(N-1)f(0,0)=\alpha$, so $f$ is constant.
This proves injectivity on the quotient by constants.
\end{proof}

The theorem constructs a global ordinary trace action and
ideal membership, not a unique expression of every defect
in moment generators. The strict order argument supplies
finite expressions after choices. Compatible primitives
in a resolution must still satisfy their higher equations.
The scalar rank $J_N(1)=N$ also remains present before
passing to the quotient by constants.

The parameter-adic completion has an exact description.
Writing $J=I_{N,\hbar}\cap W_{N,\hbar}^{GL_N}$,
saturation gives
$J\cap\hbar^mW_{N,\hbar}^{GL_N}=\hbar^mJ$.
The $R$-free quotient makes reduction of the defining
short exact sequence modulo $\hbar^m$ exact.
The transition maps $J/\hbar^{m+1}J\to J/\hbar^mJ$ are
surjective, so compatible lifts corrected in these
kernels prove exactness after inverse limit.
Thus $\Phi_{N,\hbar}$ completes to an isomorphism, and
the completed kernel is the closure of $J$.
Every product at fixed $\hbar$-order uses finitely many
polynomial coefficients. Their degrees may grow with
the order.

\clearpage
\section{Relative reduction and the remaining comparison}
\label{sec:vol2-relative-reduction-comparison}

There are now two positive comparisons.
The relative rank-two dg algebra maps strictly to
its Weyl--Cartan cohomology algebra, retaining the odd
relation module and every specialization.
The ordinary reduction at arbitrary rank has its
integral trace action across the discriminant.
The maps have different sources and preserve different
data. Neither identifies the absolute BRST complex
merely by taking rational group invariants.

For rank three, an extension requires a specified
relative quantum relation complex and a map to its
derived Cartan recipient.
Its relation cohomology and possible invariant
Massey or extension classes must be computed.
A value of one inverse component would not decide
whether such a comparison exists.

The full Hamiltonian question also retains the positive
gauge ghosts. It asks for a coherent moment comparison
on the absolute rank-two quantum recipient and its
higher-rank extensions, strict on the scalar
quadratics and with the cubic primitive of
Section~\ref{sec:vol2-derived-action-bfv}.
The comparison must carry the negative relation classes,
stabilizer states and cyclic operations across the
discriminant. The coefficient-plane parameter
$\varepsilon$ is still independent of matrix $\hbar$.
These requirements are additional to the strict relative
restriction and the ordinary trace theorem.
They are the operations to compare after the absolute
complex and its higher model have been constructed.
